%%%%%%%%%%%%%%%%%%%%%%%%%%%%%%%%%%%%%%%%%%%%%%%%%%%%%%%%%%%
\begin{frame}[fragile]\frametitle{Objective}

To understand

			\begin{center}
			\includegraphics[width=0.35\linewidth,keepaspectratio]{transformer}
			
			{\tiny From ``Attention is all you need'' paper by Vaswani, et al., 2017}
			\end{center}		

			Thats' it \ldots
			
\end{frame}

%%%%%%%%%%%%%%%%%%%%%%%%%%%%%%%%%%%%%%%%%%%%%%%%%%%%%%%%%%%
\begin{frame}[fragile]\frametitle{Assuming Familiarity with \ldots}

\begin{itemize}
\item Linear Algebra: Vectors, Matrices
\item Calculus: Gradient Descent
\item Statistics: Distributions
\item Machine Learning: Supervised, Softmax
\item Deep Learning: Back-propagation, 
\item Natural Language Processing: Tokenization, Word Vectors, Seq2Seq (RNN, LSTM)
\end{itemize}	

\end{frame}


%%%%%%%%%%%%%%%%%%%%%%%%%%%%%%%%%%%%%%%%%%%%%%%%%%%%%%%%%%%
\begin{frame}[fragile]\frametitle{Sequence-to-Sequence (seq2seq)}

\begin{center}
See any issues with this traditional seq2seq paradigm?
\end{center}	

\end{frame}

%%%%%%%%%%%%%%%%%%%%%%%%%%%%%%%%%%%%%%%%%%%%%%%%%%%%%%%%%%%
\begin{frame}[fragile]\frametitle{Issues with recurrent models}


Linear interaction distance. %O(sequence length) steps for distant word pairs to interact means:

\begin{itemize}
\item Hard to learn long-distance dependencies (because of vanishing gradient problems!)
\item Linear order of words is ``baked in''; not necessarily the  right way to think about sentences \ldots Meaning sentence structure of one language may not be correspondingly same as order in the other language.
\end{itemize}	 

\begin{center}
\includegraphics[width=\linewidth,keepaspectratio]{bert35}
\end{center}	

 
% {\tiny (Ref: Language \& Machine Learning - John Hewitt)}
\end{frame}

%%%%%%%%%%%%%%%%%%%%%%%%%%%%%%%%%%%%%%%%%%%%%%%%%%%%%%%%%%%
\begin{frame}[fragile]\frametitle{Issues with recurrent models}

Lack of parallelizability. %Forward and backward passes have O(sequence length) unparallelizable operations

\begin{itemize}
\item GPUs can perform a bunch of independent computations at once!
\item But future RNN hidden states can't be computed in full before past RNN
hidden states have been computed
\item Inhibits training on very large datasets!
\end{itemize}	 

\begin{center}
\includegraphics[width=\linewidth,keepaspectratio]{bert36}
\end{center}	

 
% {\tiny (Ref: Language \& Machine Learning - John Hewitt)}
\end{frame}

%%%%%%%%%%%%%%%%%%%%%%%%%%%%%%%%%%%%%%%%%%%%%%%%%%%%%%%%%%%
\begin{frame}[fragile]\frametitle{Then?}

If not recurrence, then what? How about word windows? Word window models aggregate local contexts

\begin{itemize}
\item Also known as 1D convolution
\item Number of unparallelizable operations not tied to sequence length!
\end{itemize}	 

\begin{center}
\includegraphics[width=\linewidth,keepaspectratio]{bert37}
\end{center}	

 
% {\tiny (Ref: Language \& Machine Learning - John Hewitt)}
\end{frame}

%%%%%%%%%%%%%%%%%%%%%%%%%%%%%%%%%%%%%%%%%%%%%%%%%%%%%%%%%%%
\begin{frame}[fragile]\frametitle{Then?}

What about long-distance dependencies?

\begin{itemize}
\item Stacking word window layers allows interaction between farther words
\item But if your sequences are too long, you'll just ignore long-distance context

\end{itemize}	 

\begin{center}
\includegraphics[width=\linewidth,keepaspectratio]{bert38}
\end{center}	

 
% {\tiny (Ref: Language \& Machine Learning - John Hewitt)}
\end{frame}

%%%%%%%%%%%%%%%%%%%%%%%%%%%%%%%%%%%%%%%%%%%%%%%%%%%%%%%%%%%
\begin{frame}[fragile]\frametitle{Attention}

If not recurrence, then what? How about attention?

\begin{itemize}
\item Attention treats each word's representation as a query to access and  incorporate information from a set of values.
\item We saw attention from the decoder to the encoder; today we'll think about
attention within a single sentence.
\item If attention gives us access to any state… maybe we can just use attention and don't need the RNN?
%\item Number of unparallelizable operations not tied to sequence length.
\item All words interact at every layer!

\end{itemize}	 

\begin{center}
\includegraphics[width=\linewidth,keepaspectratio]{bert39}
\end{center}	

 
% {\tiny (Ref: Language \& Machine Learning - John Hewitt)}
\end{frame}

%%%%%%%%%%%%%%%%%%%%%%%%%%%%%%%%%%%%%%%%%%%%%%%%%%%%%%%%%%%%%%%%%%%%%%%%%%%%%%%%%%
\begin{frame}[fragile]\frametitle{}
\begin{center}
{\Large Pre-training: What, Why, and How}
\end{center}
\end{frame}

%%%%%%%%%%%%%%%%%%%%%%%%%%%%%%%%%%%%%%%%%%%%%%%%%%%%%%%%%%%
\begin{frame}[fragile]\frametitle{What is Pre-training?}

\textbf{Pre-training} is the first large-scale training phase where the model learns general language understanding from massive text corpora --- \emph{before} any task-specific fine-tuning.

\begin{itemize}
\item The model sees billions of tokens and learns statistical patterns of language.
\item No hand-labelled task data is needed; supervision comes from the text itself (self-supervised learning).
\item The resulting weights encode broad world knowledge, grammar, reasoning patterns, etc.
\item Fine-tuning later adapts these weights cheaply to specific tasks (classification, QA, translation \ldots).
\end{itemize}

\textbf{Analogy:} Pre-training is like years of reading every book in the library. Fine-tuning is like a short specialist course on top of that foundation.

\end{frame}

%%%%%%%%%%%%%%%%%%%%%%%%%%%%%%%%%%%%%%%%%%%%%%%%%%%%%%%%%%%
\begin{frame}[fragile]\frametitle{Pre-training Objectives: Three Main Paradigms}

Different transformer architectures use different self-supervised objectives during pre-training.

\begin{center}
\begin{tabular}{|p{2.8cm}|p{3.2cm}|p{3.2cm}|}
\hline
\textbf{Objective} & \textbf{Architecture} & \textbf{Models} \\
\hline
Masked Language Modelling (MLM) & Encoder-only & BERT, RoBERTa, DistilBERT \\
\hline
Causal / Next-Token Prediction (CLM) & Decoder-only & GPT-2, GPT-3/4, LLaMA \\
\hline
Seq2Seq Denoising & Encoder-Decoder & T5, BART, PEGASUS \\
\hline
\end{tabular}
\end{center}

Each objective shapes \emph{what the model is good at} and \emph{how it should be used at inference time}.

\end{frame}

%%%%%%%%%%%%%%%%%%%%%%%%%%%%%%%%%%%%%%%%%%%%%%%%%%%%%%%%%%%
\begin{frame}[fragile]\frametitle{Pre-training: Masked Language Modelling (MLM)}

Used by \textbf{Encoder-only} models (BERT family).

\textbf{Idea:} Randomly mask 15\% of input tokens and ask the model to predict the masked words using \emph{all surrounding context} (left AND right).

\begin{center}
\texttt{Input:  "The cat [MASK] on the mat."} \\
\texttt{Target: "The cat sat  on the mat."}
\end{center}

\textbf{Why this works:}
\begin{itemize}
\item Forces the model to understand bidirectional context.
\item Every word must consider both previous and following words to predict a masked token.
\item Ideal for understanding tasks: classification, NER, extractive QA.
\end{itemize}

\textbf{Limitation:} Not naturally suited for generation --- the model was never trained to produce sequences autoregressively.

\end{frame}

%%%%%%%%%%%%%%%%%%%%%%%%%%%%%%%%%%%%%%%%%%%%%%%%%%%%%%%%%%%
\begin{frame}[fragile]\frametitle{Pre-training: Causal Language Modelling (CLM)}

Used by \textbf{Decoder-only} models (GPT family, LLaMA).

\textbf{Idea:} Predict the \emph{next} token given all \emph{previous} tokens only (left context).

\begin{center}
\texttt{Input:  "The cat sat on the"}\\
\texttt{Target: "mat"}
\end{center}

More concretely, for each position $t$ in the sequence the model predicts token $x_t$ from $x_1, x_2, \ldots, x_{t-1}$.

\textbf{Why this works:}
\begin{itemize}
\item The training signal comes from every token in every sentence (dense supervision).
\item The causal (left-to-right) mask prevents the model from ``seeing the future''.
\item Naturally suited for generation: the model learned to continue text at inference time.
\end{itemize}

\textbf{Limitation:} Only sees left context, so less optimal for pure understanding tasks.

\end{frame}

%%%%%%%%%%%%%%%%%%%%%%%%%%%%%%%%%%%%%%%%%%%%%%%%%%%%%%%%%%%
\begin{frame}[fragile]\frametitle{Pre-training: Seq2Seq Denoising (Span Corruption)}

Used by \textbf{Encoder-Decoder} models (T5, BART).

\textbf{Idea:} Corrupt the input by masking or replacing spans of text; the decoder reconstructs the original (or a cleaner version).

\textbf{T5 span corruption example:}
\begin{center}
\texttt{Input:  "The $<$X$>$ sat $<$Y$>$ mat."}\\
\texttt{Target: "$<$X$>$ cat $<$Y$>$ on the"}
\end{center}

\textbf{BART denoising example (sentence shuffling + masking):}
\begin{center}
\texttt{Input:  "mat the sat [MASK] cat on"}\\
\texttt{Target: "The cat sat on the mat."}
\end{center}

\textbf{Why this works:}
\begin{itemize}
\item Encoder learns full context understanding (bidirectional).
\item Decoder learns to generate fluent text conditioned on encoder output.
\item Best of both worlds: ideal for summarisation, translation, SQL generation.
\end{itemize}

\end{frame}

%%%%%%%%%%%%%%%%%%%%%%%%%%%%%%%%%%%%%%%%%%%%%%%%%%%%%%%%%%%
\begin{frame}[fragile]\frametitle{Pre-training Input--Output Pairs: Summary}

\begin{center}
\begin{tabular}{|p{2.5cm}|p{3.5cm}|p{3.5cm}|}
\hline
\textbf{Objective} & \textbf{Input (fed to model)} & \textbf{Output (label / target)} \\
\hline
MLM & Sentence with 15\% tokens replaced by \texttt{[MASK]} & Original token at each masked position \\
\hline
CLM & Tokens $x_1 \ldots x_{t-1}$ & Token $x_t$ (for every $t$) \\
\hline
Seq2Seq Denoising & Corrupted / masked input sentence & Clean reconstructed sentence \\
\hline
\end{tabular}
\end{center}

\vspace{0.3cm}
\textbf{Key insight:} In all three cases the label is \emph{derived from the raw text itself} --- no human annotation required. This is what makes it ``self-supervised''.

\end{frame}
